%% figures/fig-retadv.tex -- Figure for Sec. 2.4.
%% Pulled into notes.tex with %% figures/fig-retadv.tex -- Figure for Sec. 2.4.
%% Pulled into notes.tex with %% figures/fig-retadv.tex -- Figure for Sec. 2.4.
%% Pulled into notes.tex with %% figures/fig-retadv.tex -- Figure for Sec. 2.4.
%% Pulled into notes.tex with \input{figures/fig-retadv}.
%% Needs tikz with the arrows.meta library, loaded in the main preamble.
%%
%% Same world line as fig-simultaneity.tex, x = 0.18 t^2, so that the two
%% figures can be read against each other. The intersections of the light cone
%% through the field point with the world line are solved exactly, not placed
%% by eye. For a field point at (X,T) the null condition 0.18 t^2 = X -+ (T - t)
%% is a quadratic in t:
%%
%%   panel (a), X = 0.95, T = 1.6:
%%       retarded   0.18 t^2 - t + 0.65 = 0  ->  t = 0.7517, x = 0.1017
%%       advanced   0.18 t^2 + t - 2.55 = 0  ->  t = 1.9000, x = 0.6498
%%   panel (b), X = 0.52, T = 1.6:
%%       retarded   0.18 t^2 - t + 1.08 = 0  ->  t = 1.4678, x = 0.3878
%%       advanced   0.18 t^2 + t - 2.12 = 0  ->  t = 1.6374, x = 0.4826
%%
%% Both null rays are therefore at 45 degrees by construction, and the merging
%% of u and v in panel (b) is a computed fact about those roots.

\begin{figure}[htb]
\centering
\begin{tikzpicture}[>={Stealth[length=5pt]},scale=2.45,line join=round]

%% ---------------- panel (a): a field point off the world line ----------
\begin{scope}
  %% short stubs marking the light cone through x
  \draw[dashed,gray!65] (0.95,1.6) -- (1.16,1.81);
  \draw[dashed,gray!65] (0.95,1.6) -- (1.16,1.39);
  %% world line
  \draw[very thick] plot[domain=0.55:2.15,samples=70,smooth]
       ({0.18*\x*\x},\x);
  \node[right] at (0.85,2.18) {$\gamma$};
  %% the two null rays that matter
  \draw[->,thick] (0.1017,0.7517) -- (0.5259,1.1759);
  \draw[thick] (0.5259,1.1759) -- (0.95,1.6);
  \draw[->,thick] (0.6498,1.9000) -- (0.7999,1.7499);
  \draw[thick] (0.7999,1.7499) -- (0.95,1.6);
  %% events
  \fill (0.1017,0.7517) circle (1.6pt);
  \node[left,font=\footnotesize] at (0.08,0.7517) {$z(u)$};
  \fill (0.6498,1.9000) circle (1.6pt);
  \node[left,font=\footnotesize] at (0.62,1.92) {$z(v)$};
  \fill (0.95,1.6) circle (1.6pt);
  \node[right,font=\footnotesize] at (0.99,1.55) {$x$};
  %% which ray is which, set level and clear of the rays
  \node[right,font=\footnotesize] at (0.60,1.00) {retarded};
  \node[right,font=\footnotesize] at (0.78,1.94) {advanced};
  \node at (0.62,0.30) {(a)};
\end{scope}

%% ---------------- panel (b): the field point brought close --------------
\begin{scope}[xshift=4.15cm]
  \draw[dashed,gray!65] (0.52,1.6) -- (0.73,1.81);
  \draw[dashed,gray!65] (0.52,1.6) -- (0.73,1.39);
  \draw[very thick] plot[domain=0.55:2.15,samples=70,smooth]
       ({0.18*\x*\x},\x);
  \node[right] at (0.85,2.18) {$\gamma$};
  \draw[->,thick] (0.3878,1.4678) -- (0.4539,1.5339);
  \draw[thick] (0.4539,1.5339) -- (0.52,1.6);
  \draw[->,thick] (0.4826,1.6374) -- (0.5013,1.6187);
  \draw[thick] (0.5013,1.6187) -- (0.52,1.6);
  \fill (0.3878,1.4678) circle (1.6pt);
  \node[left,font=\footnotesize] at (0.36,1.42) {$z(u)$};
  \fill (0.4826,1.6374) circle (1.6pt);
  \node[left,font=\footnotesize] at (0.44,1.74) {$z(v)$};
  \fill (0.52,1.6) circle (1.6pt);
  \node[right,font=\footnotesize] at (0.76,1.55) {$x$};
  \node[align=center,font=\footnotesize] at (0.66,0.80)
       {$u$ and $v$ merge\\as $x\to\gamma$};
  \node at (0.62,0.30) {(b)};
\end{scope}

\end{tikzpicture}
\caption{What the retarded and advanced potentials are, geometrically. The light
cone through a field point $x$ meets the world line twice: once in the past, at
$z(u)$, and once in the future, at $z(v)$. The first gives $A_{\rm ret}$, the
second $A_{\rm adv}$. \emph{(a)} For $x$ well off the world line the two events
are far apart, and only the retarded one is causally admissible. \emph{(b)} As
$x$ is brought towards the world line the two roots approach each other and both
approach the same event. Both potentials diverge there, at the same rate and
with the same coefficient, because both are being sourced by the same piece of
the same world line; that is why the singular parts cancel in the difference
$\tfrac12(A_{\rm ret}-A_{\rm adv})$ and survive undiminished in the sum
$\tfrac12(A_{\rm ret}+A_{\rm adv})$. Equation~(\ref{eq:retadvsplit}) splits them
where it does for this reason and no other.
\label{fig:retadv}}
\end{figure}
.
%% Needs tikz with the arrows.meta library, loaded in the main preamble.
%%
%% Same world line as fig-simultaneity.tex, x = 0.18 t^2, so that the two
%% figures can be read against each other. The intersections of the light cone
%% through the field point with the world line are solved exactly, not placed
%% by eye. For a field point at (X,T) the null condition 0.18 t^2 = X -+ (T - t)
%% is a quadratic in t:
%%
%%   panel (a), X = 0.95, T = 1.6:
%%       retarded   0.18 t^2 - t + 0.65 = 0  ->  t = 0.7517, x = 0.1017
%%       advanced   0.18 t^2 + t - 2.55 = 0  ->  t = 1.9000, x = 0.6498
%%   panel (b), X = 0.52, T = 1.6:
%%       retarded   0.18 t^2 - t + 1.08 = 0  ->  t = 1.4678, x = 0.3878
%%       advanced   0.18 t^2 + t - 2.12 = 0  ->  t = 1.6374, x = 0.4826
%%
%% Both null rays are therefore at 45 degrees by construction, and the merging
%% of u and v in panel (b) is a computed fact about those roots.

\begin{figure}[htb]
\centering
\begin{tikzpicture}[>={Stealth[length=5pt]},scale=2.45,line join=round]

%% ---------------- panel (a): a field point off the world line ----------
\begin{scope}
  %% short stubs marking the light cone through x
  \draw[dashed,gray!65] (0.95,1.6) -- (1.16,1.81);
  \draw[dashed,gray!65] (0.95,1.6) -- (1.16,1.39);
  %% world line
  \draw[very thick] plot[domain=0.55:2.15,samples=70,smooth]
       ({0.18*\x*\x},\x);
  \node[right] at (0.85,2.18) {$\gamma$};
  %% the two null rays that matter
  \draw[->,thick] (0.1017,0.7517) -- (0.5259,1.1759);
  \draw[thick] (0.5259,1.1759) -- (0.95,1.6);
  \draw[->,thick] (0.6498,1.9000) -- (0.7999,1.7499);
  \draw[thick] (0.7999,1.7499) -- (0.95,1.6);
  %% events
  \fill (0.1017,0.7517) circle (1.6pt);
  \node[left,font=\footnotesize] at (0.08,0.7517) {$z(u)$};
  \fill (0.6498,1.9000) circle (1.6pt);
  \node[left,font=\footnotesize] at (0.62,1.92) {$z(v)$};
  \fill (0.95,1.6) circle (1.6pt);
  \node[right,font=\footnotesize] at (0.99,1.55) {$x$};
  %% which ray is which, set level and clear of the rays
  \node[right,font=\footnotesize] at (0.60,1.00) {retarded};
  \node[right,font=\footnotesize] at (0.78,1.94) {advanced};
  \node at (0.62,0.30) {(a)};
\end{scope}

%% ---------------- panel (b): the field point brought close --------------
\begin{scope}[xshift=4.15cm]
  \draw[dashed,gray!65] (0.52,1.6) -- (0.73,1.81);
  \draw[dashed,gray!65] (0.52,1.6) -- (0.73,1.39);
  \draw[very thick] plot[domain=0.55:2.15,samples=70,smooth]
       ({0.18*\x*\x},\x);
  \node[right] at (0.85,2.18) {$\gamma$};
  \draw[->,thick] (0.3878,1.4678) -- (0.4539,1.5339);
  \draw[thick] (0.4539,1.5339) -- (0.52,1.6);
  \draw[->,thick] (0.4826,1.6374) -- (0.5013,1.6187);
  \draw[thick] (0.5013,1.6187) -- (0.52,1.6);
  \fill (0.3878,1.4678) circle (1.6pt);
  \node[left,font=\footnotesize] at (0.36,1.42) {$z(u)$};
  \fill (0.4826,1.6374) circle (1.6pt);
  \node[left,font=\footnotesize] at (0.44,1.74) {$z(v)$};
  \fill (0.52,1.6) circle (1.6pt);
  \node[right,font=\footnotesize] at (0.76,1.55) {$x$};
  \node[align=center,font=\footnotesize] at (0.66,0.80)
       {$u$ and $v$ merge\\as $x\to\gamma$};
  \node at (0.62,0.30) {(b)};
\end{scope}

\end{tikzpicture}
\caption{What the retarded and advanced potentials are, geometrically. The light
cone through a field point $x$ meets the world line twice: once in the past, at
$z(u)$, and once in the future, at $z(v)$. The first gives $A_{\rm ret}$, the
second $A_{\rm adv}$. \emph{(a)} For $x$ well off the world line the two events
are far apart, and only the retarded one is causally admissible. \emph{(b)} As
$x$ is brought towards the world line the two roots approach each other and both
approach the same event. Both potentials diverge there, at the same rate and
with the same coefficient, because both are being sourced by the same piece of
the same world line; that is why the singular parts cancel in the difference
$\tfrac12(A_{\rm ret}-A_{\rm adv})$ and survive undiminished in the sum
$\tfrac12(A_{\rm ret}+A_{\rm adv})$. Equation~(\ref{eq:retadvsplit}) splits them
where it does for this reason and no other.
\label{fig:retadv}}
\end{figure}
.
%% Needs tikz with the arrows.meta library, loaded in the main preamble.
%%
%% Same world line as fig-simultaneity.tex, x = 0.18 t^2, so that the two
%% figures can be read against each other. The intersections of the light cone
%% through the field point with the world line are solved exactly, not placed
%% by eye. For a field point at (X,T) the null condition 0.18 t^2 = X -+ (T - t)
%% is a quadratic in t:
%%
%%   panel (a), X = 0.95, T = 1.6:
%%       retarded   0.18 t^2 - t + 0.65 = 0  ->  t = 0.7517, x = 0.1017
%%       advanced   0.18 t^2 + t - 2.55 = 0  ->  t = 1.9000, x = 0.6498
%%   panel (b), X = 0.52, T = 1.6:
%%       retarded   0.18 t^2 - t + 1.08 = 0  ->  t = 1.4678, x = 0.3878
%%       advanced   0.18 t^2 + t - 2.12 = 0  ->  t = 1.6374, x = 0.4826
%%
%% Both null rays are therefore at 45 degrees by construction, and the merging
%% of u and v in panel (b) is a computed fact about those roots.

\begin{figure}[htb]
\centering
\begin{tikzpicture}[>={Stealth[length=5pt]},scale=2.45,line join=round]

%% ---------------- panel (a): a field point off the world line ----------
\begin{scope}
  %% short stubs marking the light cone through x
  \draw[dashed,gray!65] (0.95,1.6) -- (1.16,1.81);
  \draw[dashed,gray!65] (0.95,1.6) -- (1.16,1.39);
  %% world line
  \draw[very thick] plot[domain=0.55:2.15,samples=70,smooth]
       ({0.18*\x*\x},\x);
  \node[right] at (0.85,2.18) {$\gamma$};
  %% the two null rays that matter
  \draw[->,thick] (0.1017,0.7517) -- (0.5259,1.1759);
  \draw[thick] (0.5259,1.1759) -- (0.95,1.6);
  \draw[->,thick] (0.6498,1.9000) -- (0.7999,1.7499);
  \draw[thick] (0.7999,1.7499) -- (0.95,1.6);
  %% events
  \fill (0.1017,0.7517) circle (1.6pt);
  \node[left,font=\footnotesize] at (0.08,0.7517) {$z(u)$};
  \fill (0.6498,1.9000) circle (1.6pt);
  \node[left,font=\footnotesize] at (0.62,1.92) {$z(v)$};
  \fill (0.95,1.6) circle (1.6pt);
  \node[right,font=\footnotesize] at (0.99,1.55) {$x$};
  %% which ray is which, set level and clear of the rays
  \node[right,font=\footnotesize] at (0.60,1.00) {retarded};
  \node[right,font=\footnotesize] at (0.78,1.94) {advanced};
  \node at (0.62,0.30) {(a)};
\end{scope}

%% ---------------- panel (b): the field point brought close --------------
\begin{scope}[xshift=4.15cm]
  \draw[dashed,gray!65] (0.52,1.6) -- (0.73,1.81);
  \draw[dashed,gray!65] (0.52,1.6) -- (0.73,1.39);
  \draw[very thick] plot[domain=0.55:2.15,samples=70,smooth]
       ({0.18*\x*\x},\x);
  \node[right] at (0.85,2.18) {$\gamma$};
  \draw[->,thick] (0.3878,1.4678) -- (0.4539,1.5339);
  \draw[thick] (0.4539,1.5339) -- (0.52,1.6);
  \draw[->,thick] (0.4826,1.6374) -- (0.5013,1.6187);
  \draw[thick] (0.5013,1.6187) -- (0.52,1.6);
  \fill (0.3878,1.4678) circle (1.6pt);
  \node[left,font=\footnotesize] at (0.36,1.42) {$z(u)$};
  \fill (0.4826,1.6374) circle (1.6pt);
  \node[left,font=\footnotesize] at (0.44,1.74) {$z(v)$};
  \fill (0.52,1.6) circle (1.6pt);
  \node[right,font=\footnotesize] at (0.76,1.55) {$x$};
  \node[align=center,font=\footnotesize] at (0.66,0.80)
       {$u$ and $v$ merge\\as $x\to\gamma$};
  \node at (0.62,0.30) {(b)};
\end{scope}

\end{tikzpicture}
\caption{What the retarded and advanced potentials are, geometrically. The light
cone through a field point $x$ meets the world line twice: once in the past, at
$z(u)$, and once in the future, at $z(v)$. The first gives $A_{\rm ret}$, the
second $A_{\rm adv}$. \emph{(a)} For $x$ well off the world line the two events
are far apart, and only the retarded one is causally admissible. \emph{(b)} As
$x$ is brought towards the world line the two roots approach each other and both
approach the same event. Both potentials diverge there, at the same rate and
with the same coefficient, because both are being sourced by the same piece of
the same world line; that is why the singular parts cancel in the difference
$\tfrac12(A_{\rm ret}-A_{\rm adv})$ and survive undiminished in the sum
$\tfrac12(A_{\rm ret}+A_{\rm adv})$. Equation~(\ref{eq:retadvsplit}) splits them
where it does for this reason and no other.
\label{fig:retadv}}
\end{figure}
.
%% Needs tikz with the arrows.meta library, loaded in the main preamble.
%%
%% Same world line as fig-simultaneity.tex, x = 0.18 t^2, so that the two
%% figures can be read against each other. The intersections of the light cone
%% through the field point with the world line are solved exactly, not placed
%% by eye. For a field point at (X,T) the null condition 0.18 t^2 = X -+ (T - t)
%% is a quadratic in t:
%%
%%   panel (a), X = 0.95, T = 1.6:
%%       retarded   0.18 t^2 - t + 0.65 = 0  ->  t = 0.7517, x = 0.1017
%%       advanced   0.18 t^2 + t - 2.55 = 0  ->  t = 1.9000, x = 0.6498
%%   panel (b), X = 0.52, T = 1.6:
%%       retarded   0.18 t^2 - t + 1.08 = 0  ->  t = 1.4678, x = 0.3878
%%       advanced   0.18 t^2 + t - 2.12 = 0  ->  t = 1.6374, x = 0.4826
%%
%% Both null rays are therefore at 45 degrees by construction, and the merging
%% of u and v in panel (b) is a computed fact about those roots.

\begin{figure}[htb]
\centering
\begin{tikzpicture}[>={Stealth[length=5pt]},scale=2.45,line join=round]

%% ---------------- panel (a): a field point off the world line ----------
\begin{scope}
  %% short stubs marking the light cone through x
  \draw[dashed,gray!65] (0.95,1.6) -- (1.16,1.81);
  \draw[dashed,gray!65] (0.95,1.6) -- (1.16,1.39);
  %% world line
  \draw[very thick] plot[domain=0.55:2.15,samples=70,smooth]
       ({0.18*\x*\x},\x);
  \node[right] at (0.85,2.18) {$\gamma$};
  %% the two null rays that matter
  \draw[->,thick] (0.1017,0.7517) -- (0.5259,1.1759);
  \draw[thick] (0.5259,1.1759) -- (0.95,1.6);
  \draw[->,thick] (0.6498,1.9000) -- (0.7999,1.7499);
  \draw[thick] (0.7999,1.7499) -- (0.95,1.6);
  %% events
  \fill (0.1017,0.7517) circle (1.6pt);
  \node[left,font=\footnotesize] at (0.08,0.7517) {$z(u)$};
  \fill (0.6498,1.9000) circle (1.6pt);
  \node[left,font=\footnotesize] at (0.62,1.92) {$z(v)$};
  \fill (0.95,1.6) circle (1.6pt);
  \node[right,font=\footnotesize] at (0.99,1.55) {$x$};
  %% which ray is which, set level and clear of the rays
  \node[right,font=\footnotesize] at (0.60,1.00) {retarded};
  \node[right,font=\footnotesize] at (0.78,1.94) {advanced};
  \node at (0.62,0.30) {(a)};
\end{scope}

%% ---------------- panel (b): the field point brought close --------------
\begin{scope}[xshift=4.15cm]
  \draw[dashed,gray!65] (0.52,1.6) -- (0.73,1.81);
  \draw[dashed,gray!65] (0.52,1.6) -- (0.73,1.39);
  \draw[very thick] plot[domain=0.55:2.15,samples=70,smooth]
       ({0.18*\x*\x},\x);
  \node[right] at (0.85,2.18) {$\gamma$};
  \draw[->,thick] (0.3878,1.4678) -- (0.4539,1.5339);
  \draw[thick] (0.4539,1.5339) -- (0.52,1.6);
  \draw[->,thick] (0.4826,1.6374) -- (0.5013,1.6187);
  \draw[thick] (0.5013,1.6187) -- (0.52,1.6);
  \fill (0.3878,1.4678) circle (1.6pt);
  \node[left,font=\footnotesize] at (0.36,1.42) {$z(u)$};
  \fill (0.4826,1.6374) circle (1.6pt);
  \node[left,font=\footnotesize] at (0.44,1.74) {$z(v)$};
  \fill (0.52,1.6) circle (1.6pt);
  \node[right,font=\footnotesize] at (0.76,1.55) {$x$};
  \node[align=center,font=\footnotesize] at (0.66,0.80)
       {$u$ and $v$ merge\\as $x\to\gamma$};
  \node at (0.62,0.30) {(b)};
\end{scope}

\end{tikzpicture}
\caption{What the retarded and advanced potentials are, geometrically. The light
cone through a field point $x$ meets the world line twice: once in the past, at
$z(u)$, and once in the future, at $z(v)$. The first gives $A_{\rm ret}$, the
second $A_{\rm adv}$. \emph{(a)} For $x$ well off the world line the two events
are far apart, and only the retarded one is causally admissible. \emph{(b)} As
$x$ is brought towards the world line the two roots approach each other and both
approach the same event. Both potentials diverge there, at the same rate and
with the same coefficient, because both are being sourced by the same piece of
the same world line; that is why the singular parts cancel in the difference
$\tfrac12(A_{\rm ret}-A_{\rm adv})$ and survive undiminished in the sum
$\tfrac12(A_{\rm ret}+A_{\rm adv})$. Equation~(\ref{eq:retadvsplit}) splits them
where it does for this reason and no other.
\label{fig:retadv}}
\end{figure}
